Помимо реализации классического пуллового, линейного и стекового аллокаторов в ходе курсовой работы была предпринята попытка разработать собственный аллокатор, являющийся комбинацией некоторых существующих идей.

Нередко в процессе разработки возникает необходимость хранить массивы не только чисел, но и более сложных данных, например, пользовательских структур или классов. Эти данные объединяет то, что они имеют несколько полей, которые в памяти обычно записаны подряд. Тем не менее, это не единственный принцип, по которому можно структурировать данные. Разработанный мною аллокатор опирается на идею колоночных баз данных "--- формата организации данных, при котором информация сохраняется и упорядочивается не по строкам, как в классических реляционных базах данных, а по столбцам, то есть отдельно по каждому полю \cite{column-databases}. Анализ различных методов аллокации памяти для структур был проведён на примере пар чисел (\texttt{pair<int, int>}), как самого простого вида структур с несколькими полями. В качестве базового аллокатора был выбран пулловый аллокатор, который сочетает достаточно высокую эффективность и возможность освобождения памяти в произвольном порядке.

Для начала была разработана новая функция для тестирования аллокаторов, работающих с парами чисел. В первую очередь эта функция была запущена на приведённом в предыдущем разделе пулловом аллокаторе. В данном случае память выделялась сначала на первый, а затем на второй элемент каждой пары. Таким образом, данные в памяти оказались записаны в порядке $x_1 y_1, x_2 y_2 \ldots$

\small
\begin{minted}{cpp}
pair<int*, int*> array[n];
for (int i = 0; i < n; i++) {
	array[i] = make_pair((int*)allocator.allocate(sizeof(int)),
						 (int*)allocator.allocate(sizeof(int)));
    if (array[i].first && array[i].second) {
		*array[i].first = 5 * i - 6;
		*array[i].second = 6 * i + 7;
	}
}
\end{minted}
\normalsize

Затем был реализован непосредственно <<колоночный аллокатор>>. Его реализация схожа с обычным пулловым аллокатором, однако в качестве полей он содержит два списка свободных блоков: \texttt{free\_x} и \texttt{free\_y}, и два пула памяти: \texttt{pool\_x} и \texttt{pool\_y}, которые аналогично выделяются с помощью функции \texttt{malloc} в конструкторе аллокатора. Аллокация и деаллокация в этом аллокаторе разделены на две функции: \texttt{(de)allocate\_x} и \texttt{(de)allocate\_y}.

\small
\inputminted{cpp}{src/column_allocator.h}
\normalsize

Скорость выделения памяти не должна сильно изменяться в зависимости от того, используется ли пулловый или колоночный аллокатор, ведь выполняемые при аллокации операции практически одинаковы. Напротив, аллокаторы продемонстрировали разные результаты в зависимости от порядка, в котором память освобождалась. Для каждого из аллокаторов были проведены две проверки: для ситуации, когда память освобождается последовательно по каждой паре, и для случая, когда сначала освобождается память, занимаемая первыми элементами в каждой паре, а затем то же самое проделывается для вторых элементов в паре.

\small
\begin{minted}{cpp}
// первый случай
for (int x : order) {
	allocator.deallocate_x(array[x].first);
	allocator.deallocate_y(array[x].second);
}
// второй случай
for (int x : order) {
	allocator.deallocate_x(array[x].first);
}
for (int x : order) {
	allocator.deallocate_y(array[x].second);
}
\end{minted}
\normalsize

Результаты исследования при $n = 500000$ приведены в таблице \ref{tab:pair_results} и в диаграмме \ref{fig:pair_results}. Анализ показал, что пулловый аллокатор выполнил поставленную задачу быстрее, когда последовательно освобождалась память для каждой пары, тогда как собственный колоночный аллокатор лучше справился с задачей освобождения памяти отдельно по каждой координате. Стандартный аллокатор также выиграл при попарном освобождении памяти, однако, как и в первом исследовании, время его работы оказалось значительно больше, чем у пользовательских аллокаторов.

Сравнение с линейным и стековым аллокатором оказалось нерелевантным, поскольку они не допускают произвольный порядок освобождения памяти. Так, в стековом аллокаторе можно реализовать освобождение памяти по столбцам, только если и выделять её отдельно по каждой координате, что маловероятно в практических задачах. Однако аллокатор, работающий с парами чисел, можно разработать и на основе одного из данных аллокаторов, если в задаче не требуется освобождение произвольного участка памяти.

\small
\begin{table}[H]
	\begin{tabular}{|c|p{6.5cm}|p{6.5cm}|}
	\hline Аллокатор & Освобождение в порядке $x_1 y_1, x_2 y_2 \ldots$ & Освобождение в порядке $x_1, x_2\ldots y_1, y_2\ldots$ \\
	\hline Пулловый & 0.078153 & 0.119514 \\
	\hline Колоночный & 0.109519 & 0.074305 \\
	\hline Стандартный & 0.281971 & 0.358355 \\
	\hline
\end{tabular}
\caption{Сравнение скорости работы аллокаторов на парах чисел}
\label{tab:pair_results}
\end{table}

\normalsize
\pgfplotsset{compat=newest}

\begin{figure}[H]
\begin{tikzpicture}
	\begin{axis}[
	ybar,
	bar width=18pt,
	width=\textwidth,
	height=7.5cm,
	ymin=0,
	ylabel={Среднее время},
		legend style={at={(0.5, 0.7)},
					  anchor=south},
	symbolic x coords={
		Пулловый,
		Колоночный,
		Стандартный},
	xtick=data,
	nodes near coords={
    	\pgfmathfloattofixed{\pgfplotspointmeta}
    	\pgfmathprintnumber[fixed,precision=4]{\pgfmathresult}
	},
	nodes near coords align={vertical},
	every node near coord/.append style={font=\small},
	enlarge x limits=0.2,
	grid=both,
	]
	\addplot coordinates {
	(Пулловый,0.078153)
	(Колоночный,0.109519)
	(Стандартный,0.281971)
	};
	\addplot+[bar shift=18pt] coordinates {
	(Пулловый,0.119514)
	(Колоночный,0.074305)
	(Стандартный,0.358355)
	};
	\legend{Освобождение по строкам, Освобождение по столбцам}
	\end{axis}
\end{tikzpicture}
\caption{Сравнение скорости работы аллокаторов на парах чисел}
\label{fig:pair_results}
\end{figure}


Несмотря на то, что исследование фокусировалось только на скорости выделения и освобождения данных, стоит упомянуть, что колоночный аллокатор может оказаться более эффективным при непосредственной работе со значениями, хранящимися в массиве. Так, например, можно предположить, что он может обеспечить более высокую производительность при чтении, сортировке или суммировании значений только по одной из координат. Далеко не всегда при итерации по структурам данные в них обрабатываются одновременно, нередко они являются в достаточной степени независимыми, и потому анализ происходит по каждому полю отдельно. Именно для таких случаев может быть применён колоночный аллокатор.
